\begin{mdframed}
    
\end{mdframed}

Escriba una introducción breve que presente los dos problemas estudiados (ordenamiento de arreglos unidimensionales y multiplicación de matrices cuadradas) y el objetivo del análisis experimental: relacionar el comportamiento práctico de los algoritmos con su complejidad teórica.

% INF-221 Tarea 1 2026-2 | Tomás San Martín | Rol 202473565-9
\section{Introducción}
La notación asintótica acota el crecimiento de un algoritmo, pero no informa
sobre constantes ocultas, jerarquía de memoria ni el efecto del tipo de
entrada. Este trabajo estudia experimentalmente cuatro algoritmos de
ordenamiento y dos de multiplicación de matrices cuadradas, midiendo tiempo y
memoria sobre 72 casos por problema (Anexo A), y contrasta lo observado con las
cotas teóricas $O(n\log n)$, $O(n^3)$ y $O(n^{\log_2 7})$. El código, los datos
y los scripts de medición y graficado se incluyen en el directorio \texttt{code/}.